\documentclass[11pt,reqno]{amsart}
\usepackage{amsmath,amssymb,graphicx,amscd}
\usepackage{mathrsfs}
\usepackage{enumerate}
% pour les S des spheres
\usepackage{mathrsfs}
\usepackage{dsfont}
\usepackage{stmaryrd}
\usepackage{wrapfig}
\usepackage{pstricks}

% ----------------------------------------------------------------
\vfuzz2pt % Don't report over-full v-boxes if over-edge is small
\hfuzz2pt % Don't report over-full h-boxes if over-edge is small

\newcommand{\prob}{{\mathbb{P}}}

\renewcommand{\S}{\mathscr{S}}  % complex unit sphere

\newcommand{\CC}{{\mathbb C}}
\newcommand{\HH}{{\mathbb H}}
\newcommand{\RR}{{\mathbb R}}
\newcommand{\QQ}{{\mathbb{Q}}}
\newcommand{\NN}{{\mathbb{N}}}
\newcommand{\ZZ}{{\mathbb{Z}}}
\renewcommand{\P}{{\mathbb{P}}}
\renewcommand{\span}{{\mathrm{span}}}
\newcommand{\sgn}{\operatorname{{sgn}}}
 \DeclareMathOperator{\Prob}{\mathbb{P}}   %probability
 \DeclareMathOperator{\E}{\mathbb{E}}      %expectation
 \DeclareMathOperator{\cov}{cov}             %logarithmic integral

 \newcommand{\I}{1\!\!1}                   %Indicator function


 \newcommand{\dd}{{\mathrm{d}}}            %differential
 \newcommand{\ii}{{\mathrm{i}}}
  \newcommand{\Id}{{\mathrm{Id}}}
\newcommand{\law}{\overset{\mbox{\rm \scriptsize law}}{=}}
\newcommand{\convlaw}{\overset{\mbox{\rm \scriptsize law}}{\longrightarrow}}


\renewcommand{\Re}{{\mathfrak{Re}}}
\renewcommand{\Im}{{\mathfrak{Im}}}

% THEOREMS -------------------------------------------------------

\newtheorem{thm}{Theorem}[section]
\newtheorem{cor}[thm]{Corollary}
\newtheorem{lem}[thm]{Lemma}
\newtheorem{prop}[thm]{Proposition}
\theoremstyle{definition}
\newtheorem{defn}[thm]{Definition}
\theoremstyle{remark}
\newtheorem*{rem}{Remark}
 \newcommand{\rd}{{\mathrm{d}}}            %differential
\newcommand{\1}{\mathds{1}}


\numberwithin{equation}{section}



\begin{document}

~\vspace{-3.5cm}
\begin{center}
{\bf Probability, homework 2, due September 16.}
\end{center}
\vspace{0.2cm}
\small

\noindent{\bf Exercise 1}. 
Let $\mathcal{A}$ be a $\sigma$-algebra, $\mathbb{P}$ a probability measure and $(A_n)_{n\geq 1}$ a sequence of events in $\mathcal{A}$ which converges to $A$. Prove that
\begin{enumerate}[(i)]
\item $A\in\mathcal{A}$;
\item  $\lim_{n\to\infty}\mathbb{P}(A_n)=\mathbb{P}(A)$.\\
\end{enumerate}

\noindent{\bf Exercise 2}. Suppose a distribution function $F$ is given by
$$
F(x)=\frac{1}{4}\mathds{1}_{[0,\infty)}(x)
+\frac{1}{2}\mathds{1}_{[1,\infty)}(x)
+\frac{1}{4}\mathds{1}_{[2,\infty)}(x)
$$
What is the probability of the following events, $(-1/2,1/2)$, $(-1/2,3/2)$, $(2/3,5/2)$, $(3,\infty)$? 
\\



\noindent{\bf Exercise 3}.  Let $\mu$ be the Lebesgue measure on $\mathbb{R}$.
Build a sequence of functions $(f_n)_{n\geq 0}$, $0\leq f_n\leq 1$, such that $\int f_n\rd \mu\to 0$ but
for any $x\in\RR$, $(f_n(x))_{n\geq 0}$ does not converge.\\




\noindent{\bf Exercise 4}. Let $X$ be a random variable in ${\rm L}^1(\Omega,\mathcal{A},\Prob)$.
Let $(A_n)_{n\geq 0}$ be a sequence of events in $\mathcal{A}$ such that $\Prob(A_N)\underset{n\to\infty}{\longrightarrow}0$. Prove that $\E(X\1_{A_n})\underset{n\to\infty}{\longrightarrow}0$.\\


\noindent{\bf Exercise 5}.
Let $(d_n)_{n\geq 0}$ be a sequence in $(0,1)$, and $K_0=[0,1]$. We define iteratively $(K_n)_{n\geq 0}$
in the following way. From $K_n$, which is the union of closed disjoint intervals, we define $K_{n+1}$
by removing from each interval of $K_n$ an open interval, centered at the middle of the previous one, with length
$d_n$ times the length of the previous one. Let $K=\cap_{n\geq 0}K_n$ ($K$ is called a Cantor set).

(a) Prove that $K$ is an uncountable compact set, with empty interior, and whose points are all accumulation points

(b) What is the Lebesgue measure of $K$?\\


\noindent{\bf Exercise 6}.
Let $X$ be a nonnegative random variable. Prove that $\E(X)<+\infty$ if and only if $\sum_{n\in\mathbb{N}}\mathbb{P}(X\geq n)<\infty$.\\


\noindent{\bf Exercise 7}. {\it Convergence in measure}.
Let $(\Omega,\mathcal{A},\mu)$ be a probability space.  and $(f_n)_{n\geq 1}, f:\Omega\to\mathbb{R}$ measurable (for the Borel $\sigma$-field on $\mathbb{R}$). We say that $(f_n)_{n\geq 1}$ converges in measure to $f$ if for any $\varepsilon>0$ we have
$$
\mathbb{\mu}(|f_n-f|>\varepsilon)\underset{n\to\infty}{\longrightarrow}0.
$$
\begin{enumerate}[(i)]
\item Show that $\int|f-f_n|\rd\mu\to 0$ implies that $f_n$ converges to $f$ in measure.  Is the reciprocal true?
\item Show that if $f_n\to f$ $\mu$-almost surely, then $f_n\to f$ in measure. Is the reciprocal true?
\item Show that if $f_n\to f$ in measure, there exists a subsequence of $(f_n)_{n\geq 1}$ which converges  $\mu$-almost surely.
\item {\it (A stronger dominated convergence theorem)} We assume that $f_n\to f$ in measure and $|f_n|\leq g$ for some integrable $g:\Omega\to \mathbb{R}$, for any $n\geq 1$.
\subitem (a) Show that $|f|\leq g$ $\mu$-a.s.
\subitem (b) Deduce that $\int |f_n-f|\rd\mu\to 0$.
\end{enumerate}


\noindent{\bf Exercise 8}.  Consider a probability space $(\Omega,\mathcal{A},\mu)$ and $(A_n)_{n}$ a sequence in $\mathcal{A}$.  Let $f:\Omega\to\mathbb{R}$ be measurable (for the Borel $\sigma$-field on $\mathbb{R}$) such that $\int_{\Omega}|\1_{A_n}-f|\rd \mu\to 0$ as $n\to\infty$. Prove that there exists $A\in\mathcal{A}$ such that $f=\1_A$ $\mu$-a.s., i.e.  $\mu(f=\1_A)=1$.\\


\noindent{\bf Exercise 9}.  Consider a probability space $(E,\mathcal{A},\mu)$ and $f_n:E\to\mathbb{R}$ measurable, $n\geq 1$. Assume $f_n\to f$ $\mu$-almost surely.  Prove that for any $\varepsilon>0$
 there exists a set $A\in\mathcal{A}$ such that $\mu(A)<\varepsilon$ and the convergence $f_n\to f$ is uniform on $A^{\rm c}$.




\end{document} 